\section*{Hoofdstuk \ref{ch:g7:chemical-reactions} --- Chemische reacties: beginstoffen en reactieproducten}

\begin{solution}{exo:g7:chemical-reactions:1}
Fysisch: ijs dat smelt, suiker die \omterm{def:g2:mixing-and-dissolving:dissolve}{oplost}. Chemisch: een lucifer die
brandt, brood dat geroosterd wordt, een spijker die roest.
\end{solution}

\begin{solution}{exo:g7:chemical-reactions:2}
De \omterm{def:g7:chemical-reactions:reactant}{beginstoffen} zijn de stoffen die verbruikt worden; de
\omterm{def:g7:chemical-reactions:reactant}{reactieproducten} zijn de nieuwe stoffen die gevormd worden.
\end{solution}

\begin{solution}{exo:g7:chemical-reactions:3}
koolstof $+$ zuurstof $\longrightarrow$ koolstofdioxide.
\end{solution}

\begin{solution}{exo:g7:chemical-reactions:4}
\omterm{def:g7:chemical-reactions:reactant}{Beginstoffen}: ijzer en zuurstof. \omterm{def:g7:chemical-reactions:reactant}{Reactieproduct}: ijzeroxide.
\end{solution}

\begin{solution}{exo:g7:chemical-reactions:5}
Het beslag toont water aan; het troebele kalkwater toont koolstofdioxide
aan.
\end{solution}

\begin{solution}{exo:g7:chemical-reactions:6}
zink $+$ zoutzuur $\longrightarrow$ zinkchloride $+$ waterstof.
\end{solution}

\begin{solution}{exo:g7:chemical-reactions:7}
Voor: 4 (in het methaanmolecuul). Na: $2 + 2 = 4$ (twee in elk
watermolecuul).
\end{solution}

\begin{solution}{exo:g7:chemical-reactions:8}
Er ontstaat geen nieuwe stof: het zout is er nog, verspreid door het
water, en je krijgt het terug door het water in te dampen. Het is een
\omterm{def:g7:chemical-reactions:physical-transformation}{fysische verandering}.
\end{solution}

\begin{solution}{exo:g7:chemical-reactions:9}
Het meeste hout heeft met zuurstof gereageerd tot \omterm{def:g7:chemical-reactions:reactant}{reactieproducten} die
gassen zijn, koolstofdioxide en waterdamp, en die verdwijnen in de
\omterm{def:g4:air-a-mixture-of-gases:air}{lucht}. Alleen de as blijft achter.
\end{solution}

\begin{solution}{exo:g7:chemical-reactions:10}
waterstof $+$ zuurstof $\longrightarrow$ water. Het \omterm{def:g7:chemical-reactions:reactant}{reactieproduct}
kleurt watervrij kopersulfaat blauw.
\end{solution}

\begin{solution}{exo:g7:chemical-reactions:11}
Voor: $2 \times 2 = 4$ waterstofatomen en 2 zuurstofatomen. Na: twee
watermoleculen bevatten $2 \times 2 = 4$ waterstofatomen en $2 \times 1 =
2$ zuurstofatomen. Dezelfde \omterm{def:g7:atoms-and-molecules:atom}{atomen}, in dezelfde aantallen.
\end{solution}

\begin{solution}{exo:g7:chemical-reactions:12}
Een zuurstofatoom kan niet uit het niets komen: de tekening heeft voor
2 zuurstofatomen en na 3. Juist is: één koolstofatoom $+$ één
zuurstofmolecuul $\longrightarrow$ één koolstofdioxidemolecuul, en er
blijft niets over.
\end{solution}

\begin{solution}{pb:g7:chemical-reactions:1}
\textbf{1.} Methaan en zuurstof.

\textbf{2.} Water.

\textbf{3.} Koolstofdioxide.

\textbf{4.} Nee: het wordt niet verbruikt, en het staat niet in het
\omterm{def:g7:chemical-reactions:word-equation}{reactieschema}.

\textbf{5.} methaan $+$ zuurstof $\longrightarrow$ koolstofdioxide $+$
water.

\textbf{6.} Chemisch: twee stoffen verdwijnen (methaan, zuurstof), en er
ontstaan twee nieuwe stoffen, die met tests aangetoond worden; en er
komen warmte en licht vrij.

\textbf{7.} \ce{CH4}, \ce{O2}, \ce{CO2}, \ce{H2O}.

\textbf{8.} 1 koolstof, 4 waterstof, $2 \times 2 = 4$ zuurstof.

\textbf{9.} 1 koolstof; $2 \times 2 = 4$ waterstof; $2 + 2 \times 1 = 4$
zuurstof. Dezelfde \omterm{def:g7:atoms-and-molecules:atom}{atomen} in dezelfde aantallen: ze zijn alleen
herschikt.

\textbf{10.} $10 \times 2 = 20$ zuurstofmoleculen.

\textbf{11.} 10 koolstofdioxidemoleculen.

\textbf{12.} $10 \times 2 = 20$ watermoleculen.
\end{solution}
